\documentclass[10pt,a4paper]{article}
\usepackage[margin=1.6cm,top=1.1cm,bottom=1.1cm]{geometry}
\usepackage{amsmath,amssymb,amsthm}
\usepackage{microtype}
\setlength{\parskip}{0pt}
\linespread{0.96}\selectfont
\usepackage{etoolbox}
\AtBeginEnvironment{proof}{\setlength{\topsep}{2pt}}
\AtBeginDocument{%
  \setlength{\abovedisplayskip}{4pt}\setlength{\belowdisplayskip}{4pt}%
  \setlength{\abovedisplayshortskip}{2pt}\setlength{\belowdisplayshortskip}{2pt}}
\makeatletter
\renewcommand{\@maketitle}{\begin{center}%
  {\large\bfseries\@title\par}\vskip 0.5em
  {\normalsize\@author\par}\end{center}\vskip 0.2em}
\renewenvironment{abstract}{\small\begin{list}{}{%
  \setlength{\leftmargin}{1.2cm}\setlength{\rightmargin}{1.2cm}%
  \setlength{\topsep}{2pt}\setlength{\partopsep}{0pt}}\item\relax}{\end{list}\vskip 0.3em}
\makeatother
\usepackage[colorlinks=true,linkcolor=blue,citecolor=blue,urlcolor=blue]{hyperref}

\DeclareMathOperator{\scl}{scl}
\DeclareMathOperator{\Aut}{Aut}
\newcommand{\nc}[1]{\langle\langle #1 \rangle\rangle}
\newcommand{\genrel}[2]{\langle #1 \mid #2 \rangle}
% Single source of truth for the two relators of the paper's Theorem 1.
%
% paper/scl-counterexample.tex \input{}s this file directly and uses
% \rWord / \rPrimeWord wherever the words appear.
%
% scl/relators.py and lean/SclCounterexample/Isomorphism.lean parse this
% same file (see their READMEs) instead of hardcoding the words, so the
% three cannot silently drift apart.
\newcommand{\rWord}{aabABabABBAbaabABBAb}
\newcommand{\rPrimeWord}{aabABabABabABBAbaBAb}
 % defines \rWord, \rPrimeWord; shared with lean/ and scl/

\theoremstyle{plain}
\newtheorem{theorem}{Theorem}
\newtheorem{lemma}[theorem]{Lemma}
\newtheorem{corollary}[theorem]{Corollary}
\theoremstyle{definition}
\newtheorem{remark}[theorem]{Remark}

\title{Isomorphic one-relator groups need not have\\ relators of equal stable commutator length}
\author{Artem Semidetnov\\[-0.2em]{\small Section de math\'ematiques, Universit\'e de Gen\`eve}}
\date{}

\pagestyle{empty}
\begin{document}

\maketitle
\thispagestyle{empty}

\begin{abstract}
Heuer and L\"oh ask whether $\genrel S r \cong \genrel {S'} {r'}$, for relators $r \in F(S)'\setminus\{e\}$ and $r'\in F(S')'\setminus\{e\}$, forces $\scl_S r = \scl_{S'} r'$. We answer this negatively with an explicit pair of words of length~$20$ in $F(a,b)'$ having $\scl=1$ and $\scl=1/2$; the isomorphism is certified by six identities between freely reduced words. As a by-product we obtain a new explicit relator for which $\|G_r\| \neq 4\cdot(\scl_S r - \tfrac12)$.
\end{abstract}

\small

Let $F=F(a,b)$ be free of rank~$2$, let $F'=[F,F]$, and write $\texttt{A}=a^{-1}$, $\texttt{B}=b^{-1}$. For $r \in F\setminus\{e\}$ put $G_r := \genrel {a,b} r = F/\nc r$, where $\nc r$ denotes normal closure. Recall that $\scl_S(w) \ge 1/2$ for all $w \in F(S)'\setminus\{e\}$ [2] and that $\scl_S$ is computable [1]. Heuer and L\"oh [3, Question~1.3(1)] ask: if $r \in F(S)'\setminus\{e\}$ and $r'\in F(S')'\setminus\{e\}$ satisfy $\genrel S r \cong \genrel{S'}{r'}$, must $\scl_S r = \scl_{S'} r'$?

\begin{theorem}\label{thm:main}
Let $S=S'=\{a,b\}$ and let
\[
  r \;=\; \texttt{\rWord}, \qquad
  r' \;=\; \texttt{\rPrimeWord}.
\]
Both are cyclically reduced of length~$20$ and lie in $F'\setminus\{e\}$, and
\[
  G_r \;\cong\; G_{r'}, \qquad\text{while}\qquad \scl_S r = 1 \;\neq\; \tfrac12 = \scl_S r'.
\]
Hence the answer to Question~1.3(1) of [3] is negative.
\end{theorem}

That $r,r'\in F'$ is immediate, both having exponent sum~$0$ in $a$ and in~$b$. Since $\scl_S$ is invariant under $\Aut(F)$, conjugation and inversion, $\scl_S r \neq \scl_S r'$ shows a~fortiori that no automorphism of~$F$ carries $r$ to a conjugate of $r'^{\pm1}$: the presentations are not Nielsen equivalent, and the isomorphism is not induced by an automorphism of~$F$.

\begin{lemma}\label{lem:cert}
Let $r \in F(a,b)$, $r' \in F(x,y)$, $p,q \in F(a,b)$ and $s,t\in F(x,y)$. Let $\varphi\colon F(x,y)\to F(a,b)$ be $x\mapsto p$, $y\mapsto q$, and $\psi \colon F(a,b) \to F(x,y)$ be $a \mapsto s$, $b \mapsto t$. If
\[
 \varphi(r') \in \nc r, \quad
 \psi(r) \in \nc{r'}, \quad
 \psi(p)x^{-1},\ \psi(q)y^{-1} \in \nc{r'}, \quad
 \varphi(s)a^{-1},\ \varphi(t)b^{-1} \in \nc r,
\]
then $\genrel{x,y}{r'} \cong \genrel{a,b}{r}$.
\end{lemma}

\begin{proof}\small
As $\nc{r'}$ is generated by conjugates of $r'^{\pm1}$, which $\varphi$ carries to conjugates of $\varphi(r')^{\pm 1}$, and $\nc r$ is normal, the first condition gives $\varphi(\nc{r'}) \subseteq \nc r$; so $\varphi$ descends to $\bar\varphi \colon \genrel{x,y}{r'} \to \genrel{a,b}{r}$, and symmetrically the second gives $\bar\psi$. The third pair says $\bar\psi\bar\varphi$ fixes $x,y$, so $\bar\psi\bar\varphi=\mathrm{id}$; the fourth says $\bar\varphi\bar\psi=\mathrm{id}$. (No Hopficity is invoked: both composites are checked.)
\end{proof}

\begin{proof}[Proof that $G_r\cong G_{r'}$]\small
Apply Lemma~\ref{lem:cert} with $p = \texttt{BAbaa}$, $q=\texttt{AABabba}$, $s = \texttt{BabABab}$, $t=\texttt{baBAb}$. Each of the six memberships is witnessed by an identity $u = \prod_i w_i c_i w_i^{-1}$ of freely reduced words in which every $c_i$ is a cyclic conjugate of $r^{\pm1}$ (for $\nc r$) or of $r'^{\pm1}$ (for $\nc{r'}$), manifestly exhibiting $u$ in the relevant normal closure. In four cases the witness is a single cyclic conjugate, $w=e$:
\begin{alignat*}{3}
\varphi(s)a^{-1} &= \texttt{ABBAbaabABabABBAbaab} &&\sim r, \qquad&
\psi(p)x^{-1} &= \texttt{BabABBAbaBAbaabABabA} \sim r',\\
\varphi(t)b^{-1} &= \texttt{AABabbaBAbaBAABabbaB} &&\sim r^{-1},&
\psi(q)y^{-1} &= \texttt{BAbaBAABabABabbaBAba} \sim r'^{-1},
\end{alignat*}
where $\sim$ denotes ``is a cyclic conjugate of''. For the remaining two, $\varphi(r')$ is a product of three cyclic conjugates of $r^{\pm1}$ with trivial conjugators,
\[
\varphi(r') = \texttt{BAbaabABabABBAbaabAB}\cdot\texttt{abABBAbaabABabABBAba}\cdot\texttt{BAABabbaBAbaBAABabba},
\]
and, with $w_1 = \texttt{BabABaabABabbaBAABabABabbaBA}$, $w_2 = \texttt{BabABaabABab}$, $w_3=\texttt{BabAB}$,
\begin{align*}
\psi(r) = {}& w_1\,\texttt{ABabABBAbaBAbaabABab}\,w_1^{-1}\cdot
          w_2\,\texttt{baBAABabABabbaBAbaBA}\,w_2^{-1}\\
      &{}\cdot w_3\,\texttt{\rPrimeWord}\,w_3^{-1},
\end{align*}
the three inner words being cyclic conjugates of $r'$, $r'^{-1}$ and $r'$ respectively. Each display is an identity of freely reduced words.
\end{proof}

The values $\scl_S r = 1$ and $\scl_S r' = 1/2$ were computed with \texttt{scallop} [4], Walker's implementation of Calegari's algorithm [1]; as calibration, the same build returns $\scl(\texttt{abAB})=1/2$, $\scl([a_1,b_1][a_2,b_2])=3/2$ and $\scl(\texttt{aaaabABAbaBAAbAB})=5/8$, reproducing Example~1.1 of [3]. Note that $\scl_S r'=1/2$ realises the Duncan--Howie gap, completing the proof of Theorem~\ref{thm:main}.

\begin{corollary}\label{cor:simvol}
For $r = \texttt{\rWord}$ one has $\|G_r\| < 2 = 4\cdot\bigl(\scl_S r - \tfrac12\bigr)$; that is, Question~1.2 of [3] has a negative answer for this~$r$.
\end{corollary}

\begin{proof}\small
By [3] the simplicial volume $\|G_r\|$ depends only on the group and not on the chosen one-relator presentation with relator in the commutator subgroup, so $\|G_r\| = \|G_{r'}\|$. The weak upper bound of [3] gives $\|G_{r'}\| < 4\cdot \scl_S r' = 2$, whereas $4\cdot(\scl_S r - \tfrac12) = 2$.
\end{proof}

\begin{remark}
Corollary~\ref{cor:simvol} compares two presentations of one group and claims nothing about~$r'$. The pair was found by a search over $\Aut(F)$-orbits of words in $F'$ of length $\le 20$, on which no result above depends; that search yields further such pairs, e.g.\ $\texttt{aabaBAAbABabAB}$ and $\texttt{aabaBAAbaBAbABaabAAB}$, with $\scl=3/4$ and $2/3$.
\end{remark}

\vskip 0.4em
{\footnotesize\setlength{\parindent}{0pt}\setlength{\hangindent}{1.4em}
\textbf{References}\par
\hangindent=1.4em [1] D.~Calegari, \emph{Stable commutator length is rational in free groups}, J.~Amer. Math. Soc. \textbf{22} (2009), 941--961.\par
\hangindent=1.4em [2] A.~J.~Duncan and J.~Howie, \emph{The genus problem for one-relator products of locally indicable groups}, Math.~Z. \textbf{208} (1991), 225--237.\par
\hangindent=1.4em [3] N.~Heuer and C.~L\"oh, \emph{Simplicial volume of one-relator groups and stable commutator length}, Algebr. Geom. Topol. \textbf{22} (2022), 1615--1661.\par
\hangindent=1.4em [4] A.~Walker, \texttt{scallop}, \url{https://github.com/aldenwalker/scallop}.\par}

\end{document}
